\documentclass{article}%
\usepackage{amsmath}
\usepackage{amsfonts}
\usepackage{amssymb}
\usepackage{graphicx}%
\setcounter{MaxMatrixCols}{30}
%TCIDATA{OutputFilter=latex2.dll}
%TCIDATA{Version=5.50.0.2960}
%TCIDATA{CSTFile=40 LaTeX article.cst}
%TCIDATA{Created=Monday, February 22, 2016 09:55:26}
%TCIDATA{LastRevised=Tuesday, June 07, 2016 13:59:36}
%TCIDATA{<META NAME="GraphicsSave" CONTENT="32">}
%TCIDATA{<META NAME="SaveForMode" CONTENT="1">}
%TCIDATA{BibliographyScheme=Manual}
%TCIDATA{<META NAME="DocumentShell" CONTENT="Standard LaTeX\Blank - Standard LaTeX Article">}
%TCIDATA{Language=American English}
%BeginMSIPreambleData
\providecommand{\U}[1]{\protect\rule{.1in}{.1in}}
%EndMSIPreambleData
\newtheorem{theorem}{Theorem}
\newtheorem{acknowledgement}[theorem]{Acknowledgement}
\newtheorem{algorithm}[theorem]{Algorithm}
\newtheorem{axiom}[theorem]{Axiom}
\newtheorem{case}[theorem]{Case}
\newtheorem{claim}[theorem]{Claim}
\newtheorem{conclusion}[theorem]{Conclusion}
\newtheorem{condition}[theorem]{Condition}
\newtheorem{conjecture}[theorem]{Conjecture}
\newtheorem{corollary}[theorem]{Corollary}
\newtheorem{criterion}[theorem]{Criterion}
\newtheorem{definition}[theorem]{Definition}
\newtheorem{example}[theorem]{Example}
\newtheorem{exercise}[theorem]{Exercise}
\newtheorem{lemma}[theorem]{Lemma}
\newtheorem{notation}[theorem]{Notation}
\newtheorem{problem}[theorem]{Problem}
\newtheorem{proposition}[theorem]{Proposition}
\newtheorem{remark}[theorem]{Remark}
\newtheorem{solution}[theorem]{Solution}
\newtheorem{summary}[theorem]{Summary}
\newenvironment{proof}[1][Proof]{\noindent\textbf{#1.} }{\ \rule{0.5em}{0.5em}}
\begin{document}

\title{Recognizing cyclic matrices and a conjecture of J.G. Thompson}
\author{John D. Dixon\\School of Mathematics and Statistics, Carleton University, \\Ottawa, Ontario, Canada K1S 5B6\\email: jdixon@math.carleton.ca}
\maketitle

\begin{abstract}
In 2006 J.G. Thompson conjectured: If $F$ is a field and $A$ is in $GL(n,F)$,
then there is a permutation matrix $P$ such that $AP$ is cyclic, that is, the
minimal polynomial of $AP$ is also its characteristic polynomial\ (open
problem 16.95 in the Kourovka Notebook). The present note provides a simple
criterion for a matrix to be cyclic and uses this to prove Thompson's conjecture.

\end{abstract}



Let $A$ be an $n\times n$ matrix over a field $F$. Then $A$ is called cyclic
if there exists a vector $v\in F^{n}$ such that $F[A]v=F^{n}$; in other words,
if we define the action of the polynomial ring $F[X]$ on $F^{n}$ by $Xu:=Au$,
then the $F[X]$-module $F^{n}$ is cyclic. An equivalent condition is that the
minimal polynomial for $A$ is of degree $n$ and hence equal to the
characteristic polynomial (such matrices were classically called
nonderogatory). It is easily proved that $A$ is cyclic if and only if in the
Jordan form for $A$ there is exactly one Jordan block with eigenvalue $\alpha$
for each of the distinct eigenvalues $\alpha$ of $A$; in particular, if $A$
has no multiple eigenvalue then $A$ is cyclic. Clearly if $A$ is cyclic then
so is any matrix similar to $A$.

The difficulty with these characterisations of cyclic matrices is that all
seem to depend on properties of powers of $A$ or knowledge of the
characteristic polynomial of $A$. The following characterisation is sometimes
easier to use. We shall illustrate this by proving a conjecture of J. G.Thompson.

We first recall some definitions. We shall use the notation $A[i_{1}%
:i_{2},j_{1}:j_{2}]$ to denote the submatrix of $A$ consisting of the entries
in rows $i$ with $i_{1}\leq i\leq i_{2}$ and columns $j$ with $j_{1}\leq j\leq
j_{2}$. We shall call an $n\times n$ matrix $A$ strongly invertible if each of
the submatrices $A[1:k,1:k]$ ($k=1,...,n$) is invertible. An $n\times n$
matrix $H=[\theta_{ij}]$ is called Hessenberg if $\theta_{ij}=0$ whenever
$i\geq j+2$; we shall call $\theta_{21},\theta_{32},...,\theta_{n,n-1}$ the
subdiagonal of $H$.

\begin{proposition}
Let $A$ be an invertible $n\times n$ matrix over a field $F$. Then the
following are equivalent:

\begin{description}
\item[(i)] $A$ is cyclic;

\item[(ii)] $A$ is similar to a Hessenberg matrix whose subdiagonal entries
are all nonzero;

\item[(iii)] $A$ is similar to a Hessenberg matrix whose subdiagonal entries
are all equal to $1$;

\item[(iv)] $A$ is similar to a matrix $B$ such that $B[2:n,1:n-1]$ is
strongly invertible.
\end{description}
\end{proposition}

\begin{proof}
((i) $\implies$ (iii)): If $A$ is cyclic then there exists $v\in F^{n}$ such
that $v_{k+1}:=A^{k}v$ $(k=0,1,...,n-1)$ is a basis of $F^{n}$. Then
$V:=[v_{1},v_{2},...,v_{n}]$ is an invertible matrix. Since $Av_{k}=v_{k+1}$
for $k=1,...,n-1$ we see that $AV=VH$ where $H$ is Hessenberg of the form
\[
\left[
\begin{array}
[c]{ccccc}%
0 & 0 & \cdots & 0 & \ast\\
1 & 0 & \cdots & 0 & \ast\\
\cdots & \cdots & \cdots & \cdots & \cdots\\
0 & 0 & \cdots & 0 & \ast\\
0 & 0 & \cdots & 1 & \ast
\end{array}
\right]  .
\]


((iii) $\implies$ (i)): If $H$ is Hessenberg with each subdiagonal entry $1$
and $e_{1}:=\left[  1,0,...,0\right]  ^{T}$ then it is easily verified that
the $k$th entry of $H^{k}e_{1}$ is $1$ and each of the following entries in
$H^{k}e_{1}$ is $0$ (for $k=0,1,...,n-1$). Hence $H^{k}e_{1}$ ($k=0,1,...,n-1$%
) is a basis for $F^{n}$ and so $H$ (and hence $A$) is cyclic.

((ii) $\iff$ (iii)): One way is trivial. On the other hand if $H$ is any
Hessenberg matrix whose subdiagonal entries are nonzero then it is easily
verified that there is a diagonal matrix $D$ such that $D^{-1}HD$ is
Hessenberg with all subdiagonal entries $1$.

((ii) $\iff$ (iv)): On way is obvious. On the other hand let $\mathcal{B}$
denote the set of all $n\times n$ matrices $B$ for which $B[2:n,1:n-1]$ is
strongly invertible. We must show that each $B\in\mathcal{B}$ is similar to a
Hessenberg matrix whose subdiagonal entries are all nonzero. Using induction
it is enough to prove the following for $k=1,2,...n-1$: if $B=[\beta_{ij}%
]\in\mathcal{B}$ and the entries $\beta_{ij}\ $equal $0$ whenever $j\leq
i-2<k$ then there exists a matrix $B^{\prime}=[\beta_{ij}^{\prime}]$ similar
to $B$ such that $B^{\prime}\in\mathcal{B}$ and the entries $\beta
_{ij}^{\prime}$ equal $0$ whenever $j\leq i-2<k+1$. Note that since $B$ and
$B^{\prime}$ lie in $\mathcal{B}$ we have $\beta_{j+1},_{j}\neq0$ for all
$j<k$ and $\beta_{j+1}^{\prime},_{j}\neq0$ for all $j<k+1$.

Thus suppose that $k\geq1$ and $B=[\beta_{ij}]\in\mathcal{B}$ and the entries
$\beta_{ij}\ $equal $0$ whenever $j\leq i-2<k$. Since $B[2:n,1:n-1]$ is
strongly invertible, $B[2:k+1,1:k]$ is both upper triangular and invertible.
Therefore $\beta_{i+1,i}\neq0$ for $i=1,...,k$ and the $(k+1)$st row of $B$
has the form $(0,...,0,\beta_{k+1,k},\ast...\ast)$. A series of elementary row
operations which consist of subtracting suitable multiples of this row from
the lower rows of $B$ will reduce all entries in the $(i,k)$-positions
($i=k+2,...,n$) to $0$ (these operations do not change the zeros in the
$(i,j)$-positions for $j\leq i-2<k$). If these are followed by the
corresponding elementary column operations we obtain a matrix $B^{\prime
}=[\beta_{ij}^{\prime}]$ similar to $B$. The column operations do not affect
any entry in the first $k+1$ columns of the matrix and so the entries
$\beta_{ij}^{\prime}$ equal $0$ whenever $j\leq i-2<k+1$. Neither do these
particular row and column operations change the determinants of the
submatrices $B[2:m+1;1:m]$ ($m=1,...,n-1$) so $B^{\prime}[2:n,1:n-1]$ is
strongly invertible. Hence $B^{\prime}\in\mathcal{B}$ and the induction step
is proved. This completes the proof of the proposition.
\end{proof}

The following appears as open problem 16.95 in [1]. It was first published in
edition No. 16 of the Kourovka Notebook (2006) and appears to have remained
open until now.

\begin{corollary}
[Conjecture of J.G. Thompson]Given any invertible $n\times n$ matrix $A$ over
a field $F$ there exists a permutation matrix $P$ such that $AP$ is cyclic.
\end{corollary}

\begin{proof}
We shall prove that there exists a permutation matrix $P$ (permuting the
columns of $A$) such that $AP\in\mathcal{B}$. Since $A\ $is not singular there
is some permutation matrix $Q$ such that the $(2,1)$th entry $AQ$ is nonzero
and this is the basis for an induction proof.

We now prove for $k=1,...,n-2$ that: if $B$ is an $n\times n$ invertible
matrix such that $B[2:k+1,1:k]$ is strongly invertible, then there exists a
permutation matrix $R_{k}$ such that $(BR_{k})[2:k+2,1:k+1]$ is strongly
invertible. Indeed since $B$ is invertible, the submatrix $B[2:k+2,1:n]$ has
rank $k+1$. Since $B[2:k+2,1:k]$ is invertible, the first $k$ columns of
$B[2:k+2,1:n]$ are linearly independent, and there must be a further column,
say the $\ell$th, which is linearly independent of these first $k$ columns.
Choose any permutation matrix $R_{k}$ which leaves the first $k$ columns of
$B$ fixed and moves the $\ell$th column into the $(k+1)$st position. Then
$BR_{k}$ has the required property. In particular, if $k=n-2$ then
$BR_{n-2}\in\mathcal{B}$.

To complete the proof take $P:=QR_{1}...R_{n-2}$ and apply the theorem.
\end{proof}

\bigskip

{\large Reference}

[1] V. D. Mazurov and E. I. Khukhro (eds.). Unsolved Problems in Group Theory:
the Kourovka Notebook (No. 18). Novosibirsk, 2014.

\qquad\qquad(available at: arXiv: 1401.0300v6 [mathGR])


\end{document}